% This file should be rename with the speaker's last name. (e.g : smith)

% Write the author names in the order you wish them to appear
% and underline the one who will give the talk.

% Only the speaker's email address is required.

\begin{abstract_online}{Parametric integrals, combinatorial identities and applications}{%
        \underline{Thierry Dana-Picard}$^{1}$, David G. Zeitoun$^{2}$}{[ndp@jct.ac.il]
        }{%
        $^1$ Department of Mathematics, Jerusalem College of Technology, Jerusalem, Israel \newline{}$^2$ Orot Israel College, Elkana, Israel \newline{}}

We propose a survey of parametric integrals (aka sequences of definite integrals) studied by undergraduates in an engineering school and pre-service teachers, in a technology-rich environment.  The papers in reference provide a few examples only. Parametric integrals are interesting both for their mathematical properties, and the numerous applicable methods, and for their importance in applied science; see [1].

Let be given an either definite or improper integral of the so-called second type $I_n=\int_a^b f_n(t) \; dt$, where $a,b \in \mathbb{R}$ and $n \in \mathbb{N}$ are given. The study of the family of integrals $I_n$ can yield the following results:
\begin{enumerate}
\item An induction formula for the sequence $(I_n)$, such as $I_{n+1}=R(n)I_n$ or $I_{n+1}=u_n+R(n)I_n$, etc., where $R(n)$ is a function of the parameter $n$; see [2,3,4,5,6].
\item A closed formula for $I_n$ as a function of the parameter $n$, often using telescoping methods which lead to factorial expressions. This is the case if $R(n)$ is a rational function. In some cases, induction connects $I_{n+2}$ and $I_n$, and the usage of double factorials may yield more compact formulas. Otherwise, the study of the convergence of a series is necessary.
\item Combinatorial identities, in the case where more than one integration method can be applied.
\item New integral presentations of classical combinatorial numbers; see [3,4,5]
\end{enumerate}

Technology contributes to the study in various ways.
\begin{enumerate}
\item A Computer Algebra System provides often an interactive tutor for integration methods. Its usage for small values of the parameter helps to find a general way to compute $I_n$ as a function of $I_{n-1}, I_{n-2},...$. With this, closed formulas can be looked for.
\item The Online Encyclopedia of Integer Sequences (\url{oeis.org}). Experiments with the CAS provide the first terms of the sequences of integrals. Using the database, candidates to describe the sequence $(I_n)$ are obtained. Determination of a closed formula is made easier.
\end{enumerate}

Specific situations may appear:
 \begin{itemize}
 \item The computation of the integral for general parameter may be performed directly by the CAS. This has been the case for $I_n=\int_0^{\pi /2 } \frac {dt}{1+tan^n(t)}$ with DERIVE (it returns immediately $\pi / 4$, an answer independent of the value of the parameter!). Other CAS had hard time with this integral. The reason is that a specific theorem is implemented there; this theorem does not appear in most textbooks and is explained in [1].
 \item If the answer is readable immediately, we are done. The answer may involve special functions. For example, if $I_n=\int_0^{\pi / 2} \sin^n t \; dt$, Maple's command returns immediately $I_n=\frac{\sqrt{\pi} \; \Gamma (\frac 12 + \frac n2)}{2 \; \Gamma (1+ \frac n2)}$, providing an incitement to learn something new, the Gamma function, as an extension of the curriculum. An example is described in [5].
 \end{itemize}

We illustrate the different cases with new examples of integrals of rational functions, trigonometric functions, etc., and examples of applications in science and engineering.




\textbf{Keywords} \newline{} Parametric integrals, combinatorics, applications

\textbf{References} \newline{}
[1] \textsc{G. Boros and V. Moll}, \emph{Irresistible Integrals: Symbolics, Analysis And Experiments In The Evaluation Of Integrals}, Cambridge University Press (2004).\newline{}
[2] \textsc{Th. Dana-Picard}, Parametric integrals and symmetries of functions, \emph{Mathematics and Computer Education}, 5-12 (2005).\newline{}
[3] \textsc{Th. Dana-Picard}, Integral presentations of Catalan numbers and Wallis formula, \emph{International Journal of Mathematical Education in Science and Technology} \textbf{42} (1), 122-128 (2011). \newline{}
[4] \textsc{Th. Dana-Picard and D.G. Zeitoun}, Sequences of definite integrals, infinite series and Stirling numbers, \emph{International Journal of Mathematical Education in Science and Technology} \textbf{43} (2), 219-230 (2012). \newline{}
[5] \textsc{Th. Dana-Picard and D.G. Zeitoun}, Parametric integrals, Wallis formula and Catalan numbers, \emph{International Journal of Mathematical Education in Science and Technology} \textbf{43} (4), 515-520 (2012).\newline{}
[6] \textsc{Th. Dana-Picard and D.G. Zeitoun}, Exploration of Parametric Integrals related to a Question of Soil Mechanics, \emph{International Journal of Mathematical Education in Science and Technology} \textbf{48} (4), 617-630 (2017).\newline{}
[7] \textsc{P. Glaister}, Factorial sums, \emph{International Journal of Mathematical Education in Science and Technology} \textbf{34} (2), 250-257 (2003).
\end{abstract_online}




